\documentclass[aspectratio=169,11pt]{beamer}
\usepackage[T1]{fontenc}
\usepackage[utf8]{inputenc}
\usepackage[spanish]{babel}
\usepackage{amsmath,amssymb}
\usepackage{booktabs}
\usepackage{tikz}
\usetikzlibrary{arrows.meta,positioning,calc,babel}

\definecolor{ink}{HTML}{1C2732}
\definecolor{teal}{HTML}{4C31FF}
\definecolor{paper}{HTML}{FFFFFF}
\colorlet{muted}{ink!70!paper}
\definecolor{amber}{HTML}{CCFF90}
\definecolor{berry}{HTML}{CCFF90}
\colorlet{kitchen}{amber!40!paper}
\colorlet{hall}{amber!65!paper}
\colorlet{study}{teal!8!paper}
\colorlet{lounge}{teal!16!paper}

\setbeamercolor{normal text}{fg=ink,bg=paper}
\setbeamercolor{structure}{fg=teal}
\setbeamercolor{frametitle}{fg=ink}
\setbeamercolor{title}{fg=ink}
\setbeamercolor{subtitle}{fg=teal}
\setbeamercolor{block title}{fg=white,bg=teal}
\setbeamercolor{block body}{fg=ink,bg=teal!8!white}
\setbeamercolor{block title alerted}{fg=ink,bg=berry}
\setbeamercolor{block body alerted}{fg=ink,bg=berry!35!paper}
\setbeamercolor{block title example}{fg=white,bg=muted}
\setbeamercolor{block body example}{fg=ink,bg=muted!7!white}
\setbeamerfont{frametitle}{size=\Large,series=\bfseries}
\setbeamerfont{title}{size=\LARGE,series=\bfseries}
\setbeamertemplate{navigation symbols}{}
\setbeamertemplate{itemize items}[circle]
\setbeamertemplate{footline}{%
  \leavevmode\hbox{%
    \begin{beamercolorbox}[wd=.79\paperwidth,ht=3ex,dp=1.5ex,leftskip=1em]{normal text}
      \scriptsize\color{muted}Murdoku%
      \ifnum\value{section}>0\relax
        \enspace/\enspace\insertsectionhead
      \fi
      \ifnum\value{subsection}>0\relax
        \enspace/\enspace\insertsubsectionhead
      \fi
    \end{beamercolorbox}%
    \begin{beamercolorbox}[wd=.21\paperwidth,ht=3ex,dp=1.5ex,rightskip=1em plus1fil]{normal text}
      \hfill\scriptsize\insertframenumber\,/\,\inserttotalframenumber
    \end{beamercolorbox}}}
\setbeamersize{text margin left=8mm,text margin right=8mm}
\setlength{\emergencystretch}{1em}

\tikzset{
  stage/.style={draw=teal,fill=white,rounded corners=3pt,
    minimum height=12mm,text width=26mm,align=center,font=\small},
  flow/.style={-{Stealth[length=2mm]},thick,draw=teal},
  furniture/.style={draw=teal,fill=white!75!teal,rounded corners=2pt,
    font=\tiny,inner sep=2pt},
  person/.style={circle,fill=ink,text=white,font=\tiny\bfseries,
    inner sep=2pt,minimum size=4.5mm},
}

\newcommand{\examplegrid}{%
  \draw[step=1,ink!35] (0,0) grid (4,4);
  \draw[ink,thick] (0,0) rectangle (4,4);
}

% 0: tablero vacio de personas; 1: objetivo; 2: alternativa ilustrativa.
\newcommand{\murdokuboard}[1]{%
  \begin{tikzpicture}[x=.85cm,y=-.85cm]
    \fill[kitchen] (0,0) rectangle (2,2);
    \fill[hall] (2,0) rectangle (4,2);
    \fill[study] (0,2) rectangle (2,4);
    \fill[lounge] (2,2) rectangle (4,4);
    \draw[step=1,ink!18,thin] (0,0) grid (4,4);
    \draw[ink,very thick] (0,0) rectangle (4,4);
    \draw[ink,very thick] (2,0) -- (2,4);
    \draw[ink,very thick] (0,2) -- (4,2);
    \foreach \i in {0,1,2,3}{
      \node[font=\scriptsize,text=muted] at (\i+.5,-.3) {\i};
      \node[font=\scriptsize,text=muted] at (-.3,\i+.5) {\i};
    }
    \node[anchor=north west,font=\tiny,text=muted] at (.05,.04) {Cocina};
    \node[anchor=north west,font=\tiny,text=muted] at (2.05,.04) {Entrada};
    \node[anchor=north west,font=\tiny,text=muted] at (.05,2.04) {Estudio};
    \node[anchor=north west,font=\tiny,text=muted] at (2.05,2.04) {Sal\'on};
    \node[furniture,draw=amber,dashed,fill=hall] at (1.5,.45) {Mesa};
    \node[furniture] at (1.5,1.4) {Silla};
    \draw[teal,fill=white!75!teal,rounded corners=2pt]
      (2.1,2.25) rectangle (3.9,2.7);
    \node[font=\tiny] at (3,2.45) {Sof\'a};
    \node[furniture,draw=amber,dashed,fill=hall] at (.5,3.4) {Armario};
    \ifnum#1=1
      \node[person] at (.65,.72) {A};
      \node[person] at (1.65,1.72) {B};
    \fi
    \ifnum#1=2
      \node[person] at (1.65,1.72) {A};
      \node[person] at (.65,.72) {B};
    \fi
    \ifnum#1>0
      \node[person] at (2.65,2.78) {C};
      \node[person,fill=berry,text=ink] at (3.65,3.72) {D};
    \fi
  \end{tikzpicture}%
}

\title{Murdoku Creator}
\subtitle{Parte 2: Selecci\'on de pistas y demostraci\'on de unicidad}
\date{}

\begin{document}

\begin{frame}[plain]
  \vspace{8mm}
  {\small\color{teal}\bfseries OPTIMIZACIÓN}
  \vspace{4mm}
  \titlepage

\end{frame}

\begin{frame}{\'Indice}
  \small
  \tableofcontents[sectionstyle=show/show,subsectionstyle=show/show/show]
\end{frame}

\section{Construcci\'on del problema}

\begin{frame}{El maestro decide qu\'e pistas entregar}
  Sea \(\mathcal R\) el pool y \(m=|\mathcal R|\).
  \[
    x_r\in\{0,1\},\qquad
    x_r=1\ \Longleftrightarrow\ \text{seleccionar la pista }r
  \]
  \[
    D(x)=\sum_{r\in\mathcal R}d_rx_r,\qquad
    C(x)=\sum_{r\in\mathcal R}x_r
  \]
  \begin{block}{Cobertura narrativa}
    Cada persona distinta de la v\'ictima debe ser sujeto de al menos una pista:
    \[
      \sum_{\substack{r\in\mathcal R\\p(r)=p}}x_r\geq1
      \qquad \forall p\neq v.
    \]
  \end{block}
  \small El maestro inicial a\'un no sabe qu\'e selecciones son \'unicas.
\end{frame}

\begin{frame}{Acercarse a la dificultad y despu\'es reducir pistas}
  \[
    e=|D(x)-T|
  \]
  \[
    \boxed{\min\ (m+1)e+C(x)}
  \]
  \begin{itemize}
    \setlength{\itemsep}{3mm}
    \item \textbf{Primero:} minimizar el error respecto a \(T\).
    \item \textbf{Despu\'es:} minimizar el n\'umero de pistas.
    \item Como \(C(x)\leq m\), una mejora de una unidad en \(e\)
      domina cualquier cambio en \(C(x)\).
  \end{itemize}
  \begin{alertblock}{No es una igualdad obligatoria}
    Puede que no exista un puzzle \'unico con dificultad exactamente \(T\).
    Adem\'as, los pesos son una medida dise\~nada, no dificultad humana validada.
  \end{alertblock}
\end{frame}

\begin{frame}{Ejemplo de la prioridad entre dos objetivos}
  \begin{center}
    Objetivo ilustrativo: \(T=15\), pool de \(m=20\) pistas.
    \medskip

    \begin{tabular}{lrrrr}
      \toprule
      Selecci\'on & \(D(x)\) & \(e\) & \(C(x)\) & \(21e+C(x)\)\\
      \midrule
      A & 14 & 1 & 4 & 25\\
      B & 15 & 0 & 7 & 7\\
      \textbf{C} & \textbf{15} & \textbf{0} & \textbf{5} & \textbf{5}\\
      \bottomrule
    \end{tabular}
  \end{center}
  \vspace{4mm}
  \begin{itemize}
    \item B mejora a A: el menor error importa m\'as que usar pocas pistas.
    \item C mejora a B: con el mismo error, entrega menos pistas.
    \item Todav\'ia falta comprobar la \textbf{unicidad} de cada candidata.
  \end{itemize}
\end{frame}

\begin{frame}{El checker intenta encontrar otra respuesta}
  \begin{columns}[T,onlytextwidth]
    \begin{column}{.48\textwidth}
      \begin{block}{Se conserva}
        \begin{itemize}
          \item El mismo mapa y los mismos muebles.
          \item Las mismas personas.
          \item Las reglas base del checker.
          \item Todas las pistas seleccionadas.
        \end{itemize}
      \end{block}
    \end{column}
    \begin{column}{.48\textwidth}
      \begin{exampleblock}{Se proh\'ibe solo \(S^\star\)}
        Si \(q_{p,c_p^\star}\) indica que \(p\) conserva su celda objetivo:
        \[
          \sum_p q_{p,c_p^\star}\leq |P|-1.
        \]
        Al menos una persona debe cambiar de posici\'on.
      \end{exampleblock}
    \end{column}
  \end{columns}
  \vspace{4mm}
  \centering
  \textbf{Hay alternativa:} las pistas no bastan.\qquad
  \textbf{No hay alternativa:} soluci\'on \'unica.
\end{frame}

\begin{frame}{Un contraejemplo explica qu\'e informaci\'on falta}
  \begin{columns}[T,onlytextwidth]
    \begin{column}{.38\textwidth}
      \centering
      \murdokuboard{2}
      \par\smallskip
      \small Alternativa ilustrativa:\\
      Alicia y Bruno intercambian posiciones.
    \end{column}
    \begin{column}{.58\textwidth}
      \small
      Siguen siendo ciertas las pistas seleccionadas:
      \begin{itemize}
        \item Alicia est\'a en la cocina.
        \item Bruno est\'a en la cocina.
        \item Carla est\'a en la columna 2.
      \end{itemize}
      \begin{block}{Pistas del pool que la rompen}
        \(r_4\): Alicia est\'a en la fila 0.\\
        \(r_5\): Bruno est\'a sobre una silla.
        \[
          x_{r_4}+x_{r_5}\geq1
        \]
      \end{block}
      No se obliga a elegir ambas, ni se afirma que este corte
      elimine todas las alternativas.
    \end{column}
  \end{columns}
\end{frame}

\begin{frame}{Benders l\'ogico: aprender de cada alternativa}
  \[
    B_k=\{r\in\mathcal R:\varphi_r(S_{\mathrm{alt}}^{(k)})=0\}
  \]
  \[
    \boxed{\sum_{r\in B_k}x_r\geq1}
  \]
  \begin{enumerate}
    \setlength{\itemsep}{2mm}
    \item Resolver el maestro hasta demostrar su \'optimo.
    \item Buscar una alternativa que cumpla las pistas seleccionadas.
    \item Si aparece, evaluar \textbf{todo el pool} sobre ella.
    \item Exigir alguna pista que la descarte y volver al maestro.
    \item Si no hay alternativa, terminar.
  \end{enumerate}
  \small Las pistas anteriores no tienen por qu\'e conservarse.
  Si \(B_k=\varnothing\), el pool completo es insuficiente y se informa de un error.
\end{frame}

\begin{frame}{Por qu\'e el resultado final es \'optimo}
  \begin{enumerate}
    \setlength{\itemsep}{4mm}
    \item Cada alternativa encontrada tiene que ser descartada por
      \textbf{cualquier} selecci\'on que haga \'unico a \(S^\star\).
    \item Cada corte expresa una condici\'on necesaria para conseguirlo.
    \item El maestro mantiene una relajaci\'on y se resuelve a optimalidad.
    \item Si su \'optimo ya no admite alternativas, es tambi\'en
      \'optimo entre las selecciones \'unicas del pool.
  \end{enumerate}
  \begin{block}{Garant\'ia precisa}
    Unicidad y m\'inimo lexicogr\'afico de
    \(\bigl(|D(x)-T|,C(x)\bigr)\), con mapa, personas y pool fijos.
  \end{block}
\end{frame}

\begin{frame}{L\'imites que conviene contar al presentar el resultado}
  \begin{itemize}
    \setlength{\itemsep}{4mm}
    \item El checker \textbf{no incorpora autom\'aticamente} la condici\'on
      de los dos ocupantes de la habitaci\'on de la v\'ictima.
      La unicidad se demuestra con su modelo base y las pistas seleccionadas.
    \item La frase HTML sobre la v\'ictima es informativa:
      no suma dificultad ni a\~nade una restricci\'on al checker.
    \item La fase espacial acepta una soluci\'on factible;
      el maestro requiere un \'optimo demostrado.
    \item La selecci\'on no tiene un l\'imite global de tiempo o iteraciones.
      Un estado inconcluso no se presenta como unicidad.
  \end{itemize}
\end{frame}

\begin{frame}{Qu\'e recibe finalmente el jugador}
  \begin{columns}[T,onlytextwidth]
    \begin{column}{.48\textwidth}
      \begin{block}{Un \'unico \texttt{puzzle.html}}
        \begin{itemize}
          \item Tablero y pistas seleccionadas.
          \item Pistas a la izquierda; leyendas debajo.
          \item Posiciones inicialmente ocultas.
          \item Bot\'on para mostrar u ocultar la soluci\'on.
          \item Uso sin conexi\'on.
        \end{itemize}
      \end{block}
    \end{column}
    \begin{column}{.48\textwidth}
      \begin{exampleblock}{Registro de la ejecuci\'on}
        \begin{itemize}
          \item Configuraci\'on y semilla efectivas.
          \item Soluci\'on objetivo y muebles.
          \item Pool y subconjunto elegido.
          \item Dificultad, error e iteraciones.
          \item Tiempos y estado de los solvers.
        \end{itemize}
      \end{exampleblock}
    \end{column}
  \end{columns}
  \small La soluci\'on est\'a integrada en el HTML: el bot\'on
  evita verla por accidente, pero no protege el c\'odigo fuente.
\end{frame}

\begin{frame}{Cuatro ideas para llevarse}
  \begin{enumerate}
    \setlength{\itemsep}{5mm}
    \item \textbf{Murdoku:} deducir posiciones a partir de reglas espaciales.
    \item \textbf{Enfoque inverso:} crear la respuesta antes que las pistas.
    \item \textbf{Creaci\'on:} construir un tablero factible y un pool verdadero.
    \item \textbf{Construcci\'on del problema:} seleccionar, buscar contraejemplos
      y refinar hasta demostrar unicidad y optimalidad.
  \end{enumerate}
  \vspace{4mm}
  \centering
  \large\color{teal}No basta con generar pistas correctas:\\
  deben dejar una sola respuesta.
\end{frame}

\section{Anexo t\'ecnico}

\begin{frame}{Anexo: modelo maestro compacto}
  \small
  \[
    \begin{aligned}
      \min\quad &(m+1)e+\sum_{r\in\mathcal R}x_r\\
      \text{sujeto a}\quad
      &\sum_{\substack{r\in\mathcal R\\p(r)=p}}x_r\geq1
        &&\forall p\neq v,\\
      &e=\left|\sum_{r\in\mathcal R}d_rx_r-T\right|,\\
      &\sum_{r\in B_j}x_r\geq1
        &&\forall\text{ cortes aprendidos }j,\\
      &x_r\in\{0,1\},\quad e\in\mathbb Z_{\geq0}.
    \end{aligned}
  \]
  \begin{itemize}
    \item \(D_{\max}=\sum_r d_r\).
    \item Dominio finito usado para el error:
      \(e\leq\max\{T,|D_{\max}-T|\}\).
    \item \texttt{add\_abs\_equality} implementa la igualdad absoluta.
    \item Solo se acepta \texttt{OPTIMAL} del maestro.
  \end{itemize}
\end{frame}

\begin{frame}{Anexo: checker con una soluci\'on objetivo excluida}
  \small
  \(U\): celdas ocupables;\quad
  \(q_{pc}=1\): la persona \(p\) ocupa \(c\).
  \[
    \begin{aligned}
      \sum_{c\in U}q_{pc}&=1 &&\forall p,\\
      \sum_{\substack{p,c\in U\\i(c)=i}}q_{pc}&\leq1
        &&\forall\text{ fila }i,\\
      \sum_{\substack{p,c\in U\\j(c)=j}}q_{pc}&\leq1
        &&\forall\text{ columna }j,\\
      \varphi_r(q)&=1 &&\forall r\text{ seleccionado},\\
      \sum_p q_{p,c_p^\star}&\leq |P|-1.
    \end{aligned}
  \]
  Sin objetivo de optimizaci\'on: basta una alternativa.
  Antes de excluir \(S^\star\), se comprueba que sea factible.
  \par\medskip
  \textbf{Referencias:}
  \texttt{generate-puzzle.tex},
  \texttt{generate-solution.tex} y
  \texttt{verify-puzzle.tex}.
\end{frame}

\end{document}
